% 2013 重庆卷理科第 5 题：几何体的三视图。
% 正（主）视图为 10×4 的矩形；侧（左）视图为下底 8、上底 2、高 4 的等腰梯形；
% 俯视图为 10×8 的矩形，被两条与长边平行的线分成 3、2、3 三部分。
\begin{tikzpicture}[font=\small, >=Stealth, line width=0.45pt, x=0.26cm, y=0.26cm]
  % ---- 正（主）视图：10 宽、4 高 ----
  \draw (0,0) rectangle (10,4);
  \draw[<->] (-2,0) -- (-2,4);
  \draw (-2.8,0) -- (-1.2,0);
  \draw (-2.8,4) -- (-1.2,4);
  \node[left] at (-2.8,2) {4};
  \draw[<->] (0,-2) -- (10,-2);
  \draw (0,-2.8) -- (0,-1.2);
  \draw (10,-2.8) -- (10,-1.2);
  \node[below] at (5,-2.2) {10};
  \node[below] at (5,-3.8) {正（主）视图};

  % ---- 侧（左）视图：下底 8、上底 2、高 4 的等腰梯形 ----
  \begin{scope}[shift={(13.6,0)}]
    \draw (0,0) -- (8,0) -- (5,4) -- (3,4) -- cycle;
    \draw[<->] (0,-2) -- (8,-2);
    \draw (0,-2.8) -- (0,-1.2);
    \draw (8,-2.8) -- (8,-1.2);
    \node[below] at (4,-2.2) {8};
    \node[below] at (4,-3.8) {侧（左）视图};
  \end{scope}

  % ---- 俯视图：10×8 的矩形，分成 3、2、3 ----
  \begin{scope}[shift={(0,-15.4)}]
    \draw (0,0) rectangle (10,8);
    \draw (0,5) -- (10,5);
    \draw (0,3) -- (10,3);
    \draw[<->] (-2,0) -- (-2,3);
    \draw[<->] (-2,3) -- (-2,5);
    \draw[<->] (-2,5) -- (-2,8);
    \draw (-2.8,0) -- (-1.2,0);
    \draw (-2.8,3) -- (-1.2,3);
    \draw (-2.8,5) -- (-1.2,5);
    \draw (-2.8,8) -- (-1.2,8);
    \node[left] at (-2.8,1.5) {3};
    \node[left] at (-2.8,4) {2};
    \node[left] at (-2.8,6.5) {3};
    \node[below] at (5,-1.4) {俯视图};
  \end{scope}
\end{tikzpicture}
